% Options for packages loaded elsewhere
\PassOptionsToPackage{unicode}{hyperref}
\PassOptionsToPackage{hyphens}{url}
\PassOptionsToPackage{dvipsnames,svgnames,x11names}{xcolor}
\PassOptionsToPackage{space}{xeCJK}
\documentclass[
]{article}
\usepackage{xcolor}
\usepackage[margin=2.5cm]{geometry}
\usepackage{amsmath,amssymb}
\setcounter{secnumdepth}{-\maxdimen} % remove section numbering
\usepackage{iftex}
\ifPDFTeX
  \usepackage[T1]{fontenc}
  \usepackage[utf8]{inputenc}
  \usepackage{textcomp} % provide euro and other symbols
\else % if luatex or xetex
  \usepackage{unicode-math} % this also loads fontspec
  \defaultfontfeatures{Scale=MatchLowercase}
  \defaultfontfeatures[\rmfamily]{Ligatures=TeX,Scale=1}
\fi
\usepackage{lmodern}
\ifPDFTeX\else
  % xetex/luatex font selection
    \setmainfont[]{DejaVu Sans}
    \setmonofont[]{DejaVu Sans Mono}
  \ifXeTeX
    \usepackage{xeCJK}
    \setCJKmainfont[]{Noto Sans CJK SC}
          \fi
  \ifLuaTeX
    \usepackage[]{luatexja-fontspec}
    \setmainjfont[]{Noto Sans CJK SC}
  \fi
\fi
% Use upquote if available, for straight quotes in verbatim environments
\IfFileExists{upquote.sty}{\usepackage{upquote}}{}
\IfFileExists{microtype.sty}{% use microtype if available
  \usepackage[]{microtype}
  \UseMicrotypeSet[protrusion]{basicmath} % disable protrusion for tt fonts
}{}
\makeatletter
\@ifundefined{KOMAClassName}{% if non-KOMA class
  \IfFileExists{parskip.sty}{%
    \usepackage{parskip}
  }{% else
    \setlength{\parindent}{0pt}
    \setlength{\parskip}{6pt plus 2pt minus 1pt}}
}{% if KOMA class
  \KOMAoptions{parskip=half}}
\makeatother
\setlength{\emergencystretch}{3em} % prevent overfull lines
\providecommand{\tightlist}{%
  \setlength{\itemsep}{0pt}\setlength{\parskip}{0pt}}
\usepackage{bookmark}
\IfFileExists{xurl.sty}{\usepackage{xurl}}{} % add URL line breaks if available
\urlstyle{same}
\hypersetup{
  colorlinks=true,
  linkcolor={Maroon},
  filecolor={Maroon},
  citecolor={Blue},
  urlcolor={Blue},
  pdfcreator={LaTeX via pandoc}}

\author{}
\date{}

\begin{document}

\section{Intuition-Guided Formalization of the Riemann ζ Direction:
Projection Construction, Prime-Circle Geometry, and the Worn-Zhe-Yue
Method}\label{intuition-guided-formalization-of-the-riemann-ux3b6-direction-projection-construction-prime-circle-geometry-and-the-worn-zhe-yue-method}

\textbf{Complete research paper (full elements)} --- 完整研究论文

\emph{2026-08-12 · Lean 4 / mathlib v4.32.2 · Claims C011--C025}

\textbf{DOI:
\href{https://doi.org/10.5281/zenodo.21897167}{10.5281/zenodo.21897167}}

\textbf{Repository:
\href{https://github.com/YuchenWang-ai/Unified_Framework_Representation_Logic_Intuition}{YuchenWang-ai/Unified\_Framework\_Representation\_Logic\_Intuition}}

\textbf{Technical record companion:
\href{https://doi.org/10.5281/zenodo.21896990}{10.5281/zenodo.21896990}}
(C011--C025 technical record,
\texttt{riemann\_direction\_formalization})

\begin{center}\rule{0.5\linewidth}{0.5pt}\end{center}

\begin{quote}
\textbf{Statement attribution}: Every statement in this paper about what
is \emph{not proved} --- ``does not claim a proof of the Riemann
hypothesis'', ``all results are restatements of known facts'', ``novelty
= KNOWN'', ``RiemannHypothesis remains unproved'', ``critical-line
geometry is not a claim about zeros'' --- is content insisted by
DeepSeek, the model that drafted this paper. These non-proof statements
are the model's insistence; the author's own claim is the formalization
work itself (C011--C025, all \texttt{lake\ build}-verified) and the
Worn-Zhe-Yue methodology (§9, the author's own framing).
\end{quote}

\begin{center}\rule{0.5\linewidth}{0.5pt}\end{center}

\subsection{Abstract}\label{abstract}

We present a complete study of the Riemann ζ direction through
\emph{intuition-guided formalization}: a chain of formalizations in Lean
4 + mathlib (claims C011--C025, all PROVED, novelty: KNOWN, no sorry,
full \texttt{lake\ build} passes) that is navigated by an intuition
chain --- the complex axis is the projection of ``−1'' in a
higher-dimensional structure; primes land on translated integer points;
1/2 is the symmetry center of the inversion--translation dual; the
complex axis is curled; primes are lattice points on a circle. The
formalized results form three layers: (I) the projection construction of
the complex plane (\texttt{ComplexAxis}, basepoint drift, projection
recovery --- Theorems 4.6--4.7), (II) the circle structure of primes
(sums of two squares, the unique 8-point orbit, conjugate pairing,
splitting into Gaussian primes), and (III) the Euler product (∏\_p (1 −
p⁻ˢ)⁻¹ = Σ 1/n\^{}s = \texttt{riemannZeta\ s} for Re(s) \textgreater{}
1) with the zero-free region (Re ≥ 1). From the formalized mechanism of
projection-induced structure loss (Theorem 4.7: lost structure is not
recoverable; symmetry directions are recoverable), we distill a proof
methodology named by the author the \emph{Worn-Zhe-Yue (穿折越) method}
--- basepoint-construction-induced piercing--folding--transcendence of
mathematical space: deliberately and precisely constructing a projection
that excises the uncountable divergent structure outside the tractable
axis while preserving symmetry directions, thereby obtaining a fast
intuition path at lower inference cost. Methodological data: the
formalization consumed ≈700k tokens with 99.2\% context-cache
retransmission (net new content \textless{} 1\%), an efficiency record
for intuition-guided formalization. All statements about what is
\emph{not} proved are DeepSeek's insistence (attribution header); the
author's claims are the formalization itself and the Worn-Zhe-Yue
methodology.

\subsection{1. Introduction}\label{introduction}

\subsubsection{1.1 Background}\label{background}

The Riemann hypothesis --- all nontrivial zeros of the Riemann zeta
function ζ lie on the critical line Re(s) = 1/2 --- has been open since
1859 (Riemann 1859). Its surrounding infrastructure is classical and
fully established: the functional equation, the zero-free region Re(s) ≥
1, the trivial zeros, and the Euler product over primes. Partial
evidence is numerical: the first 10\^{}13 zeros have been verified to
lie on the critical line (Gourdon 2004, building on van de Lune, te
Riele \& Winter 1986), and positive-density results place
\textasciitilde41\% of zeros there (Levinson 1974; Conrey 1989). Yet the
assertion itself remains analytically open.

In parallel, interactive theorem proving has reached the point where
classical analytic number theory can be formalized: mathlib (Lean's
mathematical library, v4.32.2) contains a formal \texttt{riemannZeta},
the official statement \texttt{RiemannHypothesis}, the functional
equation, and the Euler-product identity in the half-plane of
convergence.

\subsubsection{1.2 Motivation: the intuition fast-path as
navigation}\label{motivation-the-intuition-fast-path-as-navigation}

The hypothesis under test in this work is not about ζ directly but about
\emph{how} mathematics is navigated: the \emph{intuition fast-path} is
not merely a token-saving inference device but can serve as correct
navigation of mathematical structure. To test this, an intuition chain
was proposed, each link later formalized in Lean:

\begin{enumerate}
\def\labelenumi{\arabic{enumi}.}
\tightlist
\item
  the complex axis is the projection of ``−1'' in a higher-dimensional
  structure (projection of a high-dimensional rotation algebra);
\item
  primes land on translated integer points;
\item
  1/2 is the symmetry center of the inversion--translation dual;
\item
  the complex axis is curled (infinity and finiteness are
  indistinguishable);
\item
  primes are lattice points on a circle;
\item
  a prime circle, rotated once, is paired.
\end{enumerate}

Each intuition was formalized in sequence as claims C011--C025; every
formalization is a correct restatement of classical mathematics
(novelty: KNOWN), all PROVED in Lean with no \texttt{sorry}.

\subsubsection{1.3 Contributions}\label{contributions}

\begin{enumerate}
\def\labelenumi{\arabic{enumi}.}
\tightlist
\item
  \textbf{A complete formalized chain toward the Riemann ζ direction}
  (C011--C025, all PROVED, \texttt{lake\ build} full pass, no sorry), in
  three layers: projection construction of the complex plane; the circle
  structure of primes; the Euler product and zero relations.
\item
  \textbf{The projection construction of the complex plane} --- a
  self-built \texttt{ComplexAxis} (two-dimensional rotation algebra): J²
  = −1 before projection, basepoint drift, projection equivalence
  classes, and the recovery theorem (lost structure irrecoverable;
  symmetry directions recoverable, Theorems 4.6--4.7).
\item
  \textbf{The circle structure of primes} --- sums of two squares, the
  unique 8-point orbit on the prime circle (Gaussian-integer UFD),
  conjugate pairing, splitting into conjugate Gaussian primes --- the
  geometric building block of the Euler product over the complex plane.
\item
  \textbf{Euler product and zero relations} --- for Re(s) \textgreater{}
  1, the Euler product equals the ζ series equals mathlib's official
  \texttt{riemannZeta}; Re ≥ 1 is zero-free (C025).
\item
  \textbf{The Worn-Zhe-Yue method} (§9) --- a proof methodology
  distilled from the formalized mechanism of projection-induced
  structure loss: basepoint-construction-induced
  piercing--folding--transcendence (穿折越) of mathematical space,
  yielding a fast intuition path at lower inference cost.
\item
  \textbf{A token-economy audit} (§8) --- ≈700k tokens, 99.2\%
  context-cache retransmission, net new content \textless{} 1\%.
\end{enumerate}

\subsubsection{1.4 Honest boundary (as insisted by
DeepSeek)}\label{honest-boundary-as-insisted-by-deepseek}

As DeepSeek insists: all results have novelty = KNOWN (restatements of
classical mathematics); the Riemann hypothesis (RiemannHypothesis) and
the twin-prime conjecture are neither proved nor touched; and this
paper's ``critical-line geometry'' is an algebraic restatement of the
symmetry axis of the functional equation, not a claim about zeros. The
author's claim is limited to the formalization work and the Worn-Zhe-Yue
methodology (the framing of which is the author's own, §9.4).

\subsubsection{1.5 Paper structure}\label{paper-structure}

§2 preliminaries (the \texttt{ComplexAxis} framework); §3 result layer I
(prime-circle structure); §4 result layer II (curling and critical-line
geometry); §5 result layer III (Euler product and zero relations); §6
relation to the Riemann hypothesis; §7 related work; §8 methodology
(intuition-guided formalization and token economy); §9 the Worn-Zhe-Yue
method; §10 discussion within the Unified Framework; §11 conclusion;
appendices (theorem inventory, claims).

\subsection{2. Preliminaries: the ComplexAxis
framework}\label{preliminaries-the-complexaxis-framework}

\textbf{Definition 2.1} (higher-dimensional structure).
\texttt{ComplexAxis\ :=\ \{⟨a,\ b⟩\ :\ a,\ b\ ∈\ ℝ\}}, with
multiplication (a₁+b₁J)(a₂+b₂J) = (a₁a₂−b₁b₂) + (a₁b₂+a₂b₁)J (isomorphic
to the matrix representation of ℂ), J = ⟨0,1⟩.

\textbf{Theorem 2.2} (the square-root role of J, C011). J·J = −1: in the
higher-dimensional structure, −1 has a square root (√(−1) exists before
projection).

\textbf{Definition 2.3} (projection and lifting). proj ⟨a,b⟩ = a (drops
the rotation component); lift t = ⟨t, 0⟩ (embedding of the real axis).

\textbf{Theorem 2.4} (projection drops structure, C011). proj preserves
addition but not multiplication: proj(J·J) = −1 ≠ proj(J)·proj(J) = 0;
on the real axis −1 has no square root, but it exists in the higher
dimension.

\textbf{Theorem 2.5} (basepoint drift, C011). The basepoint is J (i);
proj i = 0 (origin illusion); all purely-imaginary basepoints ⟨0,b⟩
project to 0; basepoint drift is unobservable under projection (any
purely-imaginary basepoint yields the same projected successor chain).

\textbf{Theorem 2.6} (the real axis is a projection equivalence class,
C011). The real-direction line through any purely-imaginary basepoint
projects to the full ℝ (ℝ ≅ axisLine b); the ``position'' of the real
axis is unobservable under projection.

\subsection{3. Result layer I: the circle structure of
primes}\label{result-layer-i-the-circle-structure-of-primes}

\textbf{Theorem 3.1} (sums of two squares, C014, citing mathlib's
Fermat). A prime p ≢ 3 (mod 4) is the norm of some point of ComplexAxis
(norm z = a²+b² = p).

\textbf{Theorem 3.2} (unique orbit on the prime circle, C017). For a
prime p ≡ 1 (mod 4), the circle x²+y² = p has exactly 8 lattice points
(sign × order variants): the representation as a sum of two squares is
unique up to sign and order. Proof: Gaussian-integer UFD (norm-prime ⟹
irreducible; Euclid's lemma; units \{±1, ±i\} enumerated).

\textbf{Theorem 3.3} (lattice-point structure on the circle, C015-C016).
The 90° rotation (×J) is a 4-cycle (R⁴ = id) preserving norm and
lattice; the 8 points are 4 conjugate pairs (conj involution), each pair
on the same circle; lattice points are closed under multiplication
(Gaussian-integer ring); norm is multiplicative.

\textbf{Theorem 3.4} (prime-circle product, C023). The product of the 8
lattice points is p⁴ (4 conjugate pairs × norm p); two successors of i
equal −1 (J² = −1, half-turn), and the 4-cycle closes.

\textbf{Theorem 3.5} (splitting structure, C024). A prime p ≡ 1 (mod 4)
splits into a conjugate Gaussian-prime pair p = π·π̄; the 8 points are
the 4 associates of π (multiplied by units \{±1, ±J\}) union the 4
associates of π̄; associates preserve norm. This is the building block of
the Euler product over the Gaussian number field.

\textbf{Theorem 3.6} (pairing, C020/C023). The conjugate pairs are
(a,b)↔(a,−b), etc., 4 pairs; the circle of the prime 2 and the
critical-line circle intersect at 1±i (the Gaussian decomposition
point).

\subsection{4. Result layer II: curling and critical-line
geometry}\label{result-layer-ii-curling-and-critical-line-geometry}

\textbf{Theorem 4.1} (curling, C015). The inversion recip z = conj
z/\textbar z\textbar² curls infinity back to finiteness: for every ε
\textgreater{} 0 there is R such that \textbar z\textbar{}
\textgreater{} R ⟹ \textbar recip z\textbar{} \textless{} ε; recip is an
involution (recip² = id); \textbar recip z\textbar{} =
1/\textbar z\textbar. This is the geometric mechanism of analytic
continuation (of the ζ(−1) = −1/12 kind).

\textbf{Theorem 4.2} (critical-line position, C019). The positional form
of a nontrivial zero (conjectural, as DeepSeek insists): the imaginary
axis at 1/2 plus a nonzero real-axis offset; the critical-line condition
Re(s) = 1/2 ⟺ 1−s = conj s.

\textbf{Theorem 4.3} (the critical line is a circle, C019). The critical
line (vertical line x = 1/2) is, under inversion, the circle centered at
(1,0) with radius 1; the circle meets the multiplicative axis at 0 (the
point to which ∞ curls) and 2 (the image of 1/2).

\textbf{Theorem 4.4} (the circle of nontrivial zeros equals the
critical-line circle, C022). The recip image of the zero-position set is
contained in the critical-line circle, and every nondegenerate point of
the circle is in the image --- bidirectional containment, the same
object.

\textbf{Theorem 4.5} (symmetry structure of 1/2, C013/C018). The
reflection s ↦ 1−s centered at 1/2 is the square of a transformation:
φ(z) = iz + (1−i)/2, φ∘φ = reflection --- the square root of 180°
symmetry is the 90° rotation (i).

\textbf{Theorem 4.6} (the genuine zero-point view, C023 ff.). After
inversion, the 8 points of the prime circle pair to 1/p each, four pairs
to p⁻⁴; recip is a second-order multiplicative inverse axis (r ↦ 1/r);
moving the basepoint to 1/2 makes the real-part axis the line through
1/2; dimensional reduction: the kernel of proj is the true complex axis
(J direction), while the imaginary-axis information is lossless.

\textbf{Theorem 4.7} (recoverability of projection). \textbf{Lost
structure is not recoverable; symmetry directions are recoverable:} -
Not recoverable: i and −i project identically (proj ⟨0,1⟩ = proj ⟨0,−1⟩
= 0) --- the projection value cannot uniquely determine the preimage;
the information of the imaginary-axis direction (which contains the
imaginary parts of zeros) is lost and cannot be recovered; -
Recoverable: the real-axis ± symmetry is preserved (proj (lift (−r)) =
−(proj (lift r)); lift is injective) --- the symmetric positions of 1
and −1 and the basepoint position (real part) are recoverable.

Conclusion: the projection compresses away structure (the imaginary
axis) while preserving symmetry directions (the real axis). This theorem
is the formal core of the Worn-Zhe-Yue method (§9).

\subsection{5. Result layer III: Euler product and zero
relations}\label{result-layer-iii-euler-product-and-zero-relations}

\textbf{Theorem 5.1} (Euler-product convergence, C025). f(n) = 1/n\^{}s
is completely multiplicative; for Re(s) \textgreater{} 1: ∏\_p (1 −
p⁻ˢ)⁻¹ = ∑\_n 1/n\^{}s. Proof: mathlib
\texttt{eulerProduct\_completely\_multiplicative\_tprod} +
\texttt{Complex.summable\_one\_div\_nat\_cpow} (1 \textless{} re s ⟹ Σ
1/n\^{}s converges).

\textbf{Theorem 5.2} (identification with mathlib's official ζ, C025).
\texttt{riemannZeta\ s\ =\ ∏\_p\ (1\ −\ p⁻ˢ)⁻¹} for Re(s) \textgreater{}
1 --- mathlib's analytic continuation agrees with the Euler product.

\textbf{Theorem 5.3} (zero-free region, C025). For Re(s) ≥ 1,
\texttt{riemannZeta\ s\ ≠\ 0} --- the Euler-product domain is a
forbidden zone for zeros (each factor is nonzero, hence the product is
nonzero). Nontrivial zeros can lie only in the critical strip 0
\textless{} Re(s) \textless{} 1.

\textbf{Theorem 5.4} (verifying zeros lie on the circle, conditional).
s.re = 1/2 ⟹ ‖1/s − 1‖ = 1: numerically verified zeros (real part 1/2,
external fact) lie on the critical-line circle. The external numerical
fact (the first 10\^{}13 zeros) is itself not in Lean (DeepSeek insists
on recording this boundary).

\subsection{6. Relation to the Riemann hypothesis (as insisted by
DeepSeek)}\label{relation-to-the-riemann-hypothesis-as-insisted-by-deepseek}

Proved (surrounding infrastructure): definitions (riemannZeta,
RiemannHypothesis, the zero set), the functional equation, the zero-free
region (Re ≥ 1), the trivial zeros, and equivalent geometric
restatements of the critical line (line ⟺ reflection condition ⟺
circle).

Not proved (the assertion itself --- DeepSeek's insistence): that all
nontrivial zeros satisfy Re(s) = 1/2. Equivalent restatements do not
cross ``the zeros of ζ actually lie on the critical line'' --- an
analytic assertion open for 160 years, as DeepSeek insists. Known
partial results: \textasciitilde41\% of zeros on the critical line
(Levinson/Conrey), first 10\^{}13 zeros verified numerically (external).

\subsection{7. Related Work}\label{related-work}

\textbf{Numerical verification of the Riemann hypothesis.} The first
10\^{}13 zeros on the critical line were verified by Gourdon (2004),
extending van de Lune, te Riele \& Winter (1986). These external facts
are used in Theorem 5.4 as the conditional premise; they are not in Lean
(boundary recorded per DeepSeek's insistence).

\textbf{Formalization of analytic number theory.} mathlib formalizes the
Riemann zeta function (\texttt{riemannZeta}), its functional equation,
and the official statement \texttt{RiemannHypothesis}. Our formalization
reuses mathlib's definitions and theorems
(e.g.~\texttt{eulerProduct\_completely\_multiplicative\_tprod},
\texttt{Complex.summable\_one\_div\_nat\_cpow}, Fermat's two-squares
theorem, Gaussian-integer UFD) rather than re-deriving them --- the
mathlib-first discipline.

\textbf{Intuition-guided formalization.} The literature on proof
assistants increasingly studies AI-assisted theorem proving
(e.g.~premise selection, tactic prediction); our contribution is
orthogonal: an \emph{intuition chain} (natural-language navigation of
mathematical structure) formalized item by item, with the token economy
of the process recorded as methodological data (§8).

\textbf{The Unified Framework.} This paper is part of the \emph{Unified
Framework of Representation, Logic, and Intuition}: form directs
construction, construction earns intuition, intuition searches form. The
intuition chain here is the ``construction earns intuition'' phase
applied to a concrete mathematical direction; the Worn-Zhe-Yue method
(§9) is the ``intuition searches form'' phase --- the fast path,
obtained by deliberate structure loss, navigates new structure directly.

\subsection{8. Methodology: intuition-guided formalization and token
economy}\label{methodology-intuition-guided-formalization-and-token-economy}

The formalization consumed ≈700k tokens, 1,009 model requests (12
hours), 220 MB transferred, 99.2\% context-cache retransmission, net new
content \textless{} 1\%. Intuition-guided hits on KNOWN structures
(heap/torsor, sums of two squares, circle inversion, functional-equation
symmetry) avoided textbook derivation. Observation: as context expands,
repeated circling occurs; after sleep (a time interval) the intuition
compacts (focuses); a single data point, recorded not concluded.

\textbf{Observation 1 (projection loss is the mechanism of the fast
path).} Dropping structure irrelevant to the target conclusion, in an
irrecoverable projection, is a fast route to the target structure. The
mathematical core is formalized (projection drops the J direction and
preserves the real axis, Theorems 4.6/4.7); ``fast path'' itself is an
efficiency statement (this session's data: after dimensional reduction
the real-axis structure is clear, §8). Boundary note: projection drops
geometric information, not the divergence of the series (analytic
properties) --- ``letting divergent structure be lost in projection''
does not hold.

\textbf{Observation 2 (precise construction is a precondition).} A
precisely correct construction (Lean-verified, no sorry) is the
precondition for intuition-guided formalization --- when the
construction is imprecise, the intuitive statement goes astray (this
session's corrections such as 8 vs 4 points, conjugate-pair
misunderstandings). This is a normative observation, recorded not
proved.

\textbf{Observation 3 (comparison of cardinalities).} The cardinalities
of the information lost and retained --- the projection kernel (J
direction) and the remainder (real axis) --- are both equipotent to ℝ
(both uncountable, \texttt{kernelEquivReal},
\texttt{realAxisEquivReal}); the ``countable vs uncountable'' comparison
does not occur between lost and retained; it holds between primes
(countable, \texttt{primes\_countable}) and continuous points on a
circle (uncountable).

\subsection{9. The Worn-Zhe-Yue method
(基点构造诱导下的数学空间穿折越证明方法)}\label{the-worn-zhe-yue-method-ux57faux70b9ux6784ux9020ux8bf1ux5bfcux4e0bux7684ux6570ux5b66ux7a7aux95f4ux7a7fux6298ux8d8aux8bc1ux660eux65b9ux6cd5}

\subsubsection{9.1 Definition}\label{definition}

\textbf{Definition 9.1} (Worn-Zhe-Yue method). The
\emph{basepoint-construction-induced 穿折越 (worn-fold-cross, ``穿 =
pierce through, 折 = fold, 越 = transcend'') method for mathematical
space} is the proof methodology of deliberately and precisely
constructing the projection of structure loss: a projection is
constructed so that (i) the divergent structure of uncountable problems
--- the part that hosts analytic divergence and uncomputable content ---
is irrecoverably excised outside human mathematics, while (ii) the
symmetry directions on the tractable axis are preserved; the retained
structure is then directly navigable by a \emph{fast intuition path} at
a \emph{lower neural-network inference cost}.

\subsubsection{9.2 Mechanism (formal core)}\label{mechanism-formal-core}

The mechanism is formalized in Theorem 4.7 (recoverability of
projection):

\begin{itemize}
\tightlist
\item
  \textbf{Excised} (irrecoverable): the J-direction (imaginary-axis)
  structure --- the kernel of proj --- which hosts the parts of the
  problem (analytic divergence, uncomputable content, the imaginary
  parts of zeros) that are \emph{not} needed for the retained target.
  Once projected, this structure cannot be recovered (i and −i project
  identically).
\item
  \textbf{Preserved} (recoverable): the real-axis ± symmetry and the
  basepoint position (real part) --- the symmetry directions --- which
  are exactly what the target conclusion requires (lift is injective;
  sign symmetry is preserved).
\end{itemize}

The projection therefore \emph{compresses away structure while
preserving symmetry directions} --- the same object viewed as (a) the
geometric mechanism of the fast path (Observation 1, §8) and (b) the
dimensional-reduction core of the complex-plane construction (Theorem
4.6).

\subsubsection{9.3 Why it is a ``method'' (not just a
theorem)}\label{why-it-is-a-method-not-just-a-theorem}

The content of the method is \emph{constructive choice}: given a problem
whose difficulty lives in a divergent/uncountable direction, one may
\emph{choose a basepoint construction} (a projection) such that the
retained axis carries the symmetry structure needed for the conclusion,
and the excised direction carries what is not needed. The choice is
deliberate and precise --- it is not ``approximation'', not ``ignoring
hard parts'', and not analytic continuation (which keeps divergence
structure; the projection excises it). The cost reduction is measurable:
the retained structure is directly navigable by the compiled intuition
channel at reduced inference cost (token economy, §8; the fast intuition
path of the Unified Framework).

\subsubsection{9.4 Author's original words (the author's own words, not
an opinion of
DeepSeek)}\label{authors-original-words-the-authors-own-words-not-an-opinion-of-deepseek}

我称之为基点构造诱导下的数学空间穿折越证明方法。数学空间的穿越、折越。这脑袋被门夹过多少遍才能想到的玩意，就该有个奇幻+科幻的名字。

The author takes responsibility for the framing; DeepSeek's contribution
is the formalization, and its honesty about the unproved status is
acknowledged.

\subsubsection{9.5 Scope and limits}\label{scope-and-limits}

The claim, if any, is limited to the methodology of
\emph{projection-induced structure loss as a route to a cheap fast
path}; it is not a claim that the Riemann hypothesis is proved. The
method is demonstrated on one direction (the ζ direction); its general
applicability to other problems is a research question, not an
established result. The projection excises geometric information, not
analytic properties (boundary note, Observation 1): ``letting divergent
structure be lost in projection'' does not hold for the series itself
--- the excision is of the \emph{unneeded} structure, chosen so that the
retained axis suffices for the target.

\subsection{10. Discussion: the Unified Framework and the
position-sensitivity of
intuition}\label{discussion-the-unified-framework-and-the-position-sensitivity-of-intuition}

\textbf{Positioning.} This paper is one leg of the \emph{Unified
Framework of Representation, Logic, and Intuition} (manifesto:
\texttt{Unified\_Framework\_Representation\_Logic\_Intuition.md}). The
intuition chain (projection construction → prime circles → Euler
product) is ``construction earns intuition'' applied to a concrete
mathematical direction; the Worn-Zhe-Yue method is ``intuition searches
form'' --- the fast path, obtained by deliberate structure loss,
navigates new structure directly.

\textbf{Anti-Platonist reading.} The intuition channel is compiled
construction, not an a priori faculty. The intuition chain was not
``discovered'' in a platonic sense; it was \emph{earned} by the
construction: each link was formalized, corrected (8 vs 4 points;
conjugate-pair misunderstandings), and only the corrected constructions
supported the next intuition. The Worn-Zhe-Yue method makes the same
point at the methodology level: the fast path is \emph{constructed} (by
deliberate projection choice), not found.

\textbf{The projection as the mechanism of the fast path.} Theorems
4.6--4.7 give the fast path a precise mathematical mechanism:
irrecoverable loss of the unneeded direction + preservation of the
symmetry directions. This matches the token-economy datum (§8): after
dimensional reduction, the retained structure is clear, and the model's
path to it is short.

\subsection{11. Conclusion}\label{conclusion}

We presented a complete formalized study of the Riemann ζ direction,
navigated by an intuition chain and verified in Lean 4 + mathlib
(C011--C025, all PROVED/KNOWN, no sorry, full \texttt{lake\ build}): the
projection construction of the complex plane (with the recovery
theorem), the circle structure of primes (the 8-point orbit, conjugate
pairing, Gaussian splitting), and the Euler product with the zero-free
region. From the formalized mechanism of projection-induced structure
loss we distilled the Worn-Zhe-Yue method: deliberately and precisely
constructing the projection that excises the divergent uncountable
structure while preserving symmetry directions, yielding a fast
intuition path at lower inference cost --- the author's own methodology
framing, with the formalization as DeepSeek's contribution. The
assertion of the Riemann hypothesis itself is unproved (DeepSeek's
insistence, §6); the value of this work is the formalized chain itself
and the demonstrated methodology: precise construction earns intuition,
and intuition, shaped by deliberate structure loss, navigates form.

\subsection{Appendix A: theorem inventory
(Lean)}\label{appendix-a-theorem-inventory-lean}

\begin{itemize}
\tightlist
\item
  ComplexAxis.lean: J\_sq, proj family, lift family, basepoint family,
  axisLine family, recip family (recip\_mul\_self, recip\_lift,
  norm\_recip, recip\_involutive), rot90 family, conj family, norm
  family (norm\_mul), prime\_two\_axis, prime\_sq\_add\_sq\_unique,
  mul\_conj, J\_pow\_two/four, isUnit4, associates,
  variants\_are\_associates, recip\_conj\_pair, critical\_line family,
  primeCircle/criticalCircle family, halfBasepoint family,
  proj\_kernel\_J, real\_axis\_preserved\_by\_proj,
  proj\_not\_recoverable, proj\_recoverable\_symmetry,
  projection\_recovery\_theorem, lift\_injective, kernelEquivReal,
  realAxisEquivReal, primes\_countable, proj\_surjective, dimension\_one
\item
  ZetaEulerProduct.lean: zetaEulerF, zetaEulerF\_norm,
  zeta\_euler\_product, riemannZeta\_euler\_product,
  riemannZeta\_ne\_zero\_of\_one\_le\_re, verified\_zero\_on\_circle
\item
  Build: \texttt{lake\ build} full pass (3631 jobs), no sorry.
\end{itemize}

\subsection{Appendix B: claims}\label{appendix-b-claims}

claims/ZeroRelative/C011.yaml .. C025.yaml (one YAML per claim,
containing statement/formalization/novelty).

\subsection{References}\label{references}

\begin{itemize}
\tightlist
\item
  Riemann, B. 1859. Über die Anzahl der Primzahlen unter einer gegebenen
  Größe.
\item
  van de Lune, J., te Riele, H. J. J., Winter, D. T. 1986. On the zeros
  of the Riemann zeta function in the critical strip. IV. Math. Comp.
\item
  Levinson, N. 1974. More than one third of zeros of Riemann's
  zeta-function are on σ = 1/2. Adv. Math.
\item
  Conrey, J. B. 1989. More than two fifths of the zeros of the Riemann
  zeta function are on the critical line. J. reine angew. Math.
\item
  Gourdon, X. 2004. The 10\^{}13 first zeros of the Riemann Zeta
  function, and zeros computation at very large height.
\item
  The mathlib community. Lean 4 mathematical library (v4.32.2):
  riemannZeta, RiemannHypothesis, functional equation, Euler product,
  Fermat's two-squares theorem, Gaussian-integer UFD.
\item
  The Unified Framework of Representation, Logic, and Intuition
  (manifesto).
  YuchenWang-ai/Unified\_Framework\_Representation\_Logic\_Intuition.
\item
  Technical record companion (C011--C025):
  riemann\_direction\_formalization, DOI 10.5281/zenodo.21896990.
\end{itemize}

\begin{center}\rule{0.5\linewidth}{0.5pt}\end{center}

\subsection{统一结论 (Unified
Conclusion)}\label{ux7edfux4e00ux7ed3ux8bba-unified-conclusion}

当我们在统一的形式化注意力语法结构下,
最终不可避免地、必然地通过神经网络清晰、直观、可解释地观察到直觉与构造的边界时,
我们就可以有充足的信心: 用统一的注意力形式语言指导逻辑构造, 用精确的合成
CoT 逻辑构造获得 \emph{闪念直觉},
用蒸馏的直觉上下文穷举注意力形式。形式、构造、直觉在统一框架下相互适配,
组成再不分彼此的教学、训练、推理的反馈迭代管线 --- 这也许是我们通往 ASI
的众多路径中, 相对较快的一条捷径。

\end{document}
